% !PRINTED 716
\weberpage{716}

Wenn wir die Formel (10) mit der Formel \S.\,193, (15) verbinden und die im \S.\,194 eingeführte Bezeichnung
\[
X_1 = \sum{}^n \frac{\chi_1(n)}{n}, \quad X_2 = \sum{}^n \frac{\chi_2(n)}{n}
\]
gebrauchen, so ergiebt sich
\[
g\,h = g_1 h_1 \prod X_2.
\]
Ersetzt man $g$ durch seinen Werth (14), so folgt für die Classenzahl
\weq{15}{h = k h_1,}
wenn
\weq{16}{
\begin{aligned}
k & = \frac{m\sqrt{q\varDelta_1}}{2^{\nu-1}\pi^\nu} \prod X_2 \quad \text{bei ungeradem } m \\
& = \frac{m\sqrt{\varDelta_1}}{2^{\nu-1}\pi^\nu} \prod X_2 \quad \text{bei geradem } m,
\end{aligned}
}
und worin [nach (10), (13)]:
\weq{17}{h_1 = \frac{\sqrt{\varDelta_1}}{2^{\nu-1}E} \prod X_1}
die Classenzahl des Körpers $H_m$ ist. Im letzteren Ausdrucke erstreckt sich das Product auf alle Charaktere $\chi_1$ der Gruppe $C_1$ \gesperrt{mit Ausnahme des Hauptcharakters}.

Dass $k$ eine ganze Zahl ist, und daher $h$ ein Vielfaches von $h_1$, wird sich bald zeigen. Die Zahlen $k$ und $h_1$ heissen der \gesperrt{erste und zweite Factor der Classenzahl}.

In dem Falle, dass
\[
m = 2^\varkappa
\]
eine Potenz von 2 ist, haben wir nach \S.\,176, (7):
\[
\varDelta_1 = 2^{(\varkappa-1)2^{\varkappa-2}-1},
\]
und wenn wir diesen Werth einsetzen, so können wir für die vorstehenden Formeln in diesem Falle schreiben:
\weq{18}{
\begin{aligned}
k & = \frac{\sqrt{2}\,2^{2^{\varkappa-3}(\varkappa-3)}\,2^\varkappa}{\pi^\nu} \prod X_2 \\
h_1 & = \frac{\sqrt{2}\,2^{2^{\varkappa-3}(\varkappa-3)}}{E} \prod X_1.
\end{aligned}
}
